%% ma: L12-B03-0020
%% lop: 12
%% chuong: 1
%% bai: 3
%% dang: 12.3.1
%% loai: TN4
%% muc: NB
%% dap_an: "B"
%% nguon: "Ảnh giáo viên gửi 05/11/2026 (M11): Kiểm tra giữa HK1 Toán 12 (2026-2027), THPT Hồ Thị Bi (TP.HCM), mã đề 1133, Phần I Câu 2"
%% kien_thuc: "Bài 3 (đường tiệm cận đứng đọc từ đồ thị)"
%% trang_thai: chinh-thuc
%% ghi_chu: "Đề gốc không kèm đáp án; đã giải độc lập. Hình dựng lại từ hàm số y=(2-x)/(x-1) (tiệm cận đứng x=1, tiệm cận ngang y=-1, đi qua (0;-2) và (2;0)) khớp các số liệu trên hình gốc."
\begin{ex}
Cho hàm số $y=\dfrac{ax+b}{cx+d}$ $(ac\ne0,\ ad-bc\ne0)$ có đồ thị như hình sau:
\begin{center}
\begin{tikzpicture}[scale=.8]
  \begin{scope}
    \clip (-2,-3.6) rectangle (5.2,3.5);
    \draw[thick,domain=-1.8:0.84,samples=80,smooth] plot (\x,{(2-\x)/(\x-1)});
    \draw[thick,domain=1.16:5.2,samples=80,smooth] plot (\x,{(2-\x)/(\x-1)});
  \end{scope}
  \draw[truc] (-2.1,0)--(5.3,0) node[below]{$x$};
  \draw[truc] (0,-3.7)--(0,3.8) node[left]{$y$};
  \draw[phu] (1,-3.6)--(1,3.5) (-2,-1)--(5.2,-1);
  \path (0,0) node[above left]{\footnotesize$O$} (1,0) node[below left]{\footnotesize$1$} (2,0) node[above right]{\footnotesize$2$}
        (0,-1) node[above left]{\footnotesize$-1$} (0,-2) node[below left]{\footnotesize$-2$};
\end{tikzpicture}
\end{center}
Đường tiệm cận đứng của đồ thị hàm số đã cho có phương trình là
\choice{$y=1$}{\True $x=1$}{$x=-1$}{$y=-1$}
\loigiai{Từ đồ thị, khi $x\to1$ thì $|y|\to+\infty$ nên đường thẳng $x=1$ là tiệm cận đứng.
(Đường thẳng $y=-1$ là tiệm cận ngang.)}
\end{ex}
